\begin{proposition}
\label{prop-counting-prediction-context}
%
(\textbf{Ratio of expected counts during autoregressive and bidirectional training}).
%
Let $\mathcal{V} = \{t_0, \dots, t_V\}$ be a set of tokens.
%
Let $\mathcal{U}$ be the sample space of all possible sequences of $N$ tokens, and let the sequence $U = \{t_1, \dots, t_N\} \in\mathcal{U}$ be a random variable with probability distribution $P(U)$ defined over $\mathcal{U}$.
%
Let $\Pr[t_j = t^*]$ be the probability that the token at index $j$ in a sequence $U \in \mathcal{U}$ is given by $t^* \in \mathcal{V}$.
%
Let $\mathbb{E}_{P}[\mu_c(t^*)]$ be the expected number of tokens that are predicted by a given token $t^*$, and $\mathbb{E}_{P}[\mu_p(t^*)]$ the expected number of tokens that predict a given token $t^*$.
%
For autoregressive training, the ratio
%
\begin{equation}
    \frac{\mathbb{E}_{P}[\mu_c(t^*, U)]}{\mathbb{E}_{P}[\mu_p(t^*, U)]} = \frac{\sum_{k=1}^N (N-k)\Pr[t_k = t^*]}{\sum_{k=1}^N(k-1)\Pr[t_k = t^*]} \,,
\end{equation}
%
depends on $\Pr[t_k = t^*] \,\,\forall t^* \in \mathcal{V}, \forall k \in \{1,\dots. N\}$, while for bidirectional training the same ratio is given by 
%
\begin{equation}
     \frac{\mathbb{E}_{P}[\mu_c(t^*, U)]}{\mathbb{E}_{P}[\mu_p(t^*, U)]} = 1 \,.
\end{equation}
%    
\end{proposition}

% \begin{proposition}[Ratio of expected counts]
% \label{prop:count_ratio}
% %
% Fix a vocabulary $\mathcal V=\{t_0,\dots,t_V\}$ and an integer
% sequence length $N\ge 2$.
% Let
% \(
%   (\Omega,\mathscr F,\Pr)
% \)
% be a probability space that carries  

% * random tokens
%   \(
%     T_1,\dots,T_N:\Omega\to\mathcal V,
%   \)
%   collected in the random sequence
%   \(U=(T_1,\dots,T_N)\in\mathcal V^{N}\), and  

% * (for the bidirectional‐training analysis) i.i.d.\
%   Bernoulli masking variables
%   \(
%     M_1,\dots,M_N\sim\mathrm{Bernoulli}(\rho),
%     \quad 0<\rho<1,
%   \)
%   which are independent of each other and of the tokens.

% For a fixed token
% \(t^*\in\mathcal V\) define the random variables
% \begin{align}
%   \mu_c(t^*,U)
%     &=\sum_{k=1}^{N}(N-k)\,
%        \mathbf 1_{\{T_k=t^*\}}, &\text{(tokens \emph{predicted by} }t^*),
%   \\
%   \mu_p(t^*,U)
%     &=\sum_{k=1}^{N}(k-1)\,
%        \mathbf 1_{\{T_k=t^*\}}, &\text{(tokens that \emph{predict} }t^*).
% \end{align}
% Write
% \(
%   E[\cdot]=\E_{U,M_1{:}N}[\cdot]
% \)
% for expectation with respect to all randomness.

% \begin{enumerate}[label=(\roman*),topsep=4pt,itemsep=2pt]
% \item \textbf{Autoregressive training.}
%       Then
%       \begin{equation}\label{eq:ratio_ar}
%         \frac{E[\mu_c(t^*)]}{E[\mu_p(t^*)]}
%         \;=\;
%         \frac{\displaystyle\sum_{k=1}^{N}(N-k)\Pr[T_k=t^*]}
%              {\displaystyle\sum_{k=1}^{N}(k-1)\Pr[T_k=t^*]}\;.
%       \end{equation}

% \item \textbf{Bidirectional (masked-LM) training.}
%       Assume the model predicts every \emph{masked} token using
%       every \emph{unmasked} token.
%       Under the independence assumptions stated above,
%       \[
%         \frac{E[\mu_c(t^*)]}{E[\mu_p(t^*)]}\;=\;1 .
%       \]
% \end{enumerate}
% \end{proposition}

% \begin{proof}
% \textbf{(i) Autoregressive case.}
% Because $\mathbf 1_{\{T_k=t^*\}}$ is
% $\{0,1\}$‐valued, linearity of expectation yields
% \[
%   E[\mu_c(t^*)]
%   =\sum_{k=1}^{N}(N-k)E[\mathbf 1_{\{T_k=t^*\}}]
%   =\sum_{k=1}^{N}(N-k)\Pr[T_k=t^*].
% \]
% The same calculation with the weights $(k-1)$ gives
% \(
%   \E[\mu_p(t^*)]=\sum_{k=1}^{N}(k-1)\Pr[T_k=t^*].
% \)
% Taking the ratio gives \eqref{eq:ratio_ar}.

% \medskip
% \noindent
% \textbf{(ii) Bidirectional case.}
% Fix a position $k$.
% If $T_k=t^*$ and $M_k=0$ (token visible),
% it contributes
% \(
%   \sum_{k'\ne k}\mathbf 1_{\{M_{k'}=1\}}
% \)
% to $\mu_c(t^*)$.
% If $T_k=t^*$ and $M_k=1$ (token masked),
% it contributes
% \(
%   \sum_{k'\ne k}\mathbf 1_{\{M_{k'}=0\}}
% \)
% to $\mu_p(t^*)$.
% Taking expectations and using independence of
% \(\{M_i\}\) together with
% \(E[\mathbf 1_{\{M_i=1\}}]=\rho\),
% \(E[\mathbf 1_{\{M_i=0\}}]=1-\rho\),
% we obtain
% \[
%   E[\mu_c(t^*)]
%   =\sum_{k=1}^{N}\Pr[T_k=t^*]\,(1-\rho)\,(N-1)\rho,
%   \qquad
%   E[\mu_p(t^*)]
%   =\sum_{k=1}^{N}\Pr[T_k=t^*]\,(  \rho)\,(N-1)(1-\rho).
% \]
% The multiplicative factors are identical, hence
% \(
%   E[\mu_c(t^*)]=E[\mu_p(t^*)]
% \)
% and the stated ratio equals~$1$.
% \end{proof}

% \begin{remark}[Effect of token–position statistics]\label{rem:iid_vs_corr}
% \hfill
% \begin{enumerate}[label=(\alph*),leftmargin=*,topsep=2pt,itemsep=2pt]
% \item \emph{i.i.d.\ tokens.}
%       If the tokens are independent and identically distributed,
%       \(\Pr[T_k=t^*]=\pi_{t^*}\) for all $k$.
%       Then the two sums in \eqref{eq:ratio_ar} coincide,
%       giving
%       \(
%         E[\mu_c(t^*)]=E[\mu_p(t^*)]=
%         \pi_{t^*}\tfrac{N(N-1)}2
%       \)
%       and the ratio is~1.  Hence no directional bias appears.

% \item \emph{Correlated tokens.}
%       For a general joint law
%       \(\Pr[T_1{:}T_N]\),
%       \eqref{eq:ratio_ar} remains valid but cannot be simplified.
%       Writing the joint law via the reverse chain rule
%       \(
%         \Pr[T_1{:}T_N]
%         =\prod_{i=1}^{N}\Pr[T_i\mid T_{i+1{:}N}],
%       \)
%       one obtains
%       \[
%         \Pr[T_k=t^*]
%         =\!\!\!\sum_{\substack{t_{-k}\in\mathcal V^{N-1}}}
%           \Bigl(
%             \Pr[T_k=t^*\mid T_{k+1{:}N}=t_{k+1{:}N}]
%             \!\!\!\prod_{i\ne k}
%             \Pr[T_i=t_i\mid T_{i+1{:}N}=t_{i+1{:}N}]
%           \Bigr),
%       \]
%       which may depend strongly on $k$.
%       If $t^*$ is more likely to occur at early positions
%       than late ones, the numerator of \eqref{eq:ratio_ar}
%       (weighted by $N-k$) becomes larger than the denominator,
%       giving a ratio \(>1\) and
%       inducing directionality even under autoregressive training.
% \end{enumerate}
% \end{remark}
